% 附录：核心代码（关键片段，完整代码见 code/ 目录）

\textbf{1. 孕周解析与达标指标（utils.py 核心）}

\begin{lstlisting}[language=Python]
def parse_wk(s):
    t = str(s).lower()
    if "w" not in t:
        return float("nan")
    w = float(t.split("w")[0])
    d = float(t.split("+")[1]) if "+" in t else 0.0
    return w + d / 7.0

THETA = 0.04  # Y 染色体浓度达标阈值
male["孕周"] = male["检测孕周"].apply(parse_wk)
male["达标"] = (male["Y染色体浓度"] >= THETA).astype(int)
\end{lstlisting}

\textbf{2. 问题一面板分解与混合效应模型}

\begin{lstlisting}[language=Python]
import statsmodels.formula.api as smf
# 组内去均值（固定效应）
df["y_w"] = df["y"] - df.groupby("code")["y"].transform("mean")
df["week_w"] = df["week"] - df.groupby("code")["week"].transform("mean")
fe = smf.ols("y_w ~ week_w - 1", data=df).fit()   # 组内孕周斜率
# 组间 BMI
gb = df.drop_duplicates("code")
between = smf.ols("y_mean ~ bmi_mean", data=gb).fit()  # 组间 BMI 效应
# 混合效应模型
me = smf.mixedlm("y ~ week + bmi", data=df, groups=df["code"]).fit(reml=True)
\end{lstlisting}

\textbf{3. 问题二达标时间估计与最佳时点}

\begin{lstlisting}[language=Python]
BETA1 = 0.00319  # 问题一固定效应斜率
df["adj"] = df["y"] - BETA1 * df["week"]
person = df.groupby("code").agg(a=("adj", "mean"), bmi=("bmi", "mean")).reset_index()
person["tstar"] = (THETA - person["a"]) / BETA1   # 个体达标时间
# BMI 分组，取组内 90% 分位作为最佳时点
person["group"] = pd.cut(person["bmi"], bins=[20,28,32,36,40,60], right=False)
best = person.groupby("group")["tstar"].quantile(0.9)
\end{lstlisting}

\textbf{4. 问题四女胎异常判定（随机森林 + 交叉验证）}

\begin{lstlisting}[language=Python]
from sklearn.ensemble import RandomForestClassifier
from sklearn.model_selection import StratifiedKFold, cross_val_predict
skf = StratifiedKFold(n_splits=5, shuffle=True, random_state=42)
rf = RandomForestClassifier(n_estimators=500, class_weight="balanced", random_state=42)
rf_prob = cross_val_predict(rf, X, y, cv=skf, method="predict_proba")[:, 1]
auc = roc_auc_score(y, rf_prob)   # AUC = 0.789
\end{lstlisting}
